\subsection{Measurability with respect to an integral of point measures}

PROPOSITION 3. --- Let \((\pi, g)\) be a \(\mu\) -adapted pair, and let

\[
\nu = \int g (t) \varepsilon_ {\pi (t)} d \mu (t).
\]

Let \(f\) be a mapping of \(X\) into a topological space \(G\) , and let \(S\) be the \((\mu\text{-measurable})\) set of points \(t \in T\) such that \(g(t) > 0\) . In order that \(f\) be \(\nu\) -measurable, it is necessary and sufficient that the restriction of \(f \circ \pi\) to \(S\) be \(\mu\) -measurable.

Suppose first that \(f\) is \(\nu\) -measurable. By hypothesis, the set \(\mathcal{K}\) of compact subsets \(K\) of \(S\) , such that the restriction of \(\pi\) to \(K\) is continuous, is \(\mu\) -dense in \(S\) (Ch. IV, §5, No.~10, Prop. 15). To show that the restriction of \(f \circ \pi\) to \(S\) is \(\mu\) -measurable, it therefore suffices to prove that for every \(K \in \mathcal{K}\) , the set of compact subsets \(H\) of \(K\) , such that the restriction of \(f \circ \pi\) to \(H\) is continuous, is \(\mu\) -dense in \(K\) (Ch. IV, §5, No.~8, Prop. 13). But by hypothesis, there exists a partition of the compact set \(\pi(K)\) formed by a \(\nu\) -negligible set \(N\) and a sequence of compact sets \((C_n)\) such that the restriction of \(f\) to each \(C_n\) is continuous. Under these conditions, \(K \cap \overline{\pi}^{-1}(N)\) and the sets \(K \cap \overline{\pi}^{-1}(C_n)\) form a partition of \(K\) ; but \(K \cap \overline{\pi}^{-1}(N)\) is \(\mu\) -negligible by virtue of the Cor. of Th. 1 of No.~2, the sets \(K \cap \overline{\pi}^{-1}(C_n)\) are compact, and the restriction of \(f \circ \pi\) to each of the latter sets is continuous, which proves that the restriction of \(f \circ \pi\) to \(S\) is \(\mu\) -measurable.

Conversely, suppose this to be the case; to show that \(f\) is \(\nu\) -measurable, it suffices to prove that the set \(\mathfrak{L}\) of compact subsets \(L\) of \(X\) , such that the restriction of \(f\) to \(L\) is continuous, is \(\nu\) -dense (Ch. IV, §5, No.~10, Prop. 15). Let \(N\) be a subset of \(X\) such that \(N \cap L\) is \(\nu\) -negligible for every \(L \in \mathfrak{L}\) , and let us show that \(N\) is locally \(\nu\) -negligible. For this, we must show that \(\overline{\pi}^{-1}(N) \cap S\) is locally \(\mu\) -negligible (Cor. of Th. 1 of No.~2). Now, the set \(\mathfrak{H}\) of compact subsets \(H\) of \(S\) , such that the restrictions to \(H\) of \(\pi\) and \(f \circ \pi\) are continuous, is by hypothesis \(\mu\) -dense in \(S\) (Ch. IV, §5, No.~10, Prop. 15). It therefore suffices to prove that \(\overline{\pi}^{-1}(N) \cap H\) is \(\mu\) -negligible for every \(H \in \mathfrak{H}\) . Now, \(\pi(H)\) is compact and may be identified with the quotient space of \(H\) by the equivalence relation \(\pi(t) = \pi(t')\) , \(\pi\) being identified with the canonical mapping of \(H\) onto this quotient space (GT, I, §5, No.~2, Prop. 3). Since the restriction of \(f \circ \pi\) to \(H\) is continuous, the restriction of \(f\) to \(\pi(H)\) is therefore continuous, in other words \(\pi(H) \in \mathfrak{L}\) , consequently \(N \cap \pi(H)\) is \(\nu\) -negligible. By the Cor. of Th. 1 of No.~2, \(\overline{\pi}^{-1}(N \cap \pi(H)) \cap S\) is locally \(\mu\) -negligible; the same is therefore true of the set

\[
\mathrm{H} \cap \overline {{\pi}} ^ {- 1} (\mathrm{N}) \subset \overline {{\pi}} ^ {- 1} \bigl (\mathrm{N} \cap \pi (\mathrm{H}) \bigr) \cap \mathrm{S};
\]

but since H is compact, \(\mathrm{H} \cap \overline{\pi}^{-1}(\mathrm{N})\) is \(\mu\) -negligible, which completes the proof.

Remark. --- If f is a mapping of X into a Banach space F, it comes to the same to say that the restriction of \(f \circ \pi\) to S is \(\mu\) -measurable or to say that the function \((\mathbf{f} \circ \pi)g\) (defined on T) is \(\mu\) -measurable, since \(g\) is \(\mu\) -measurable, is not zero in S, and is zero on T - S (Ch. IV, §5, No.~10, Prop. 15).

